../cictikz/src/cictikz/data/tex/cictikz_preamble.tex
% Top-level SUN_PLL, drawn from the xschem netlist in sun_pll_sky130nm.
% The arrangement follows the source schematic, which is deliberate:
% the loop reads left to right (PFD -> CP -> filter -> buffer -> ring
% oscillator), the divider sits under the oscillator and hands CK_FB
% back to the PFD, and the housekeeping (KICK start-up, BIAS) lives at
% the bottom. The LPF hangs off VLPF because it IS a shunt - the pump
% output node is the buffer input. VDD_ROSC is red: the buffered filter
% voltage is the oscillator's supply, and that supply is the control
% node. AVDD/AVSS routing is omitted; ROSC's PWRUP_1V8 connects by
% name, as in the source.
\begin{circuitikz}[american, thick, transform shape, circuit ee IEC]
../cictikz/src/cictikz/data/tex/cictikz_lib.tex

% ---- main loop --------------------------------------------------------
\draw (0,0) rectangle (2.2,2);
\node[anchor=center] at (1.1,1) {PFD};
\draw (4.6,0) rectangle (6.6,2);
\node[anchor=center] at (5.6,1) {CP};
\draw (13.2,0) rectangle (15.4,2);
\node[anchor=center] at (14.3,1) {ROSC};

% unity-gain buffer
\draw (10.0,0.8) -- (10.0,2.0) -- (11.2,1.4) -- cycle;
\node[anchor=west] at (10.05,1.4) {$\times 1$};

% PFD inputs
\draw (-0.8,1.4) \portIn{CK\_REF} (-0.8,1.4) -- (0,1.4);

% PFD -> CP
\draw (2.2,1.4) -- (4.6,1.4);
\node[anchor=south] at (3.4,1.4) {\small $\overline{CP_{UP}}$};
\draw (2.2,0.6) -- (4.6,0.6);
\node[anchor=north] at (3.4,0.6) {\small $CP_{DOWN}$};

% CP -> buffer, with the filter shunted on VLPF
\draw (6.6,1.4) -- (10.0,1.4);
\node[anchor=south] at (7.5,1.4) {\small VLPF};
\draw (7.4,-1.8) rectangle (9.0,-0.6);
\node[anchor=center] at (8.2,-1.2) {LPF};
\draw (8.2,-0.6) to [short,-] ++(0,2.0);
\fill (8.2,1.4) circle (0.075);
\draw (6.6,0.6) -- ++(2.0,0) -- (8.6,-0.6);
\node[anchor=south] at (7.5,0.6) {\small VLPFZ};

% buffer -> oscillator: the control node
\draw (11.2,1.4) -- (13.2,1.4);
\node[anchor=south, red] at (12.2,1.4) {\small VDD\_ROSC};

% oscillator output and feedback path
\draw (15.4,1.4) -- ++(1.2,0) \portOut{CK};
\draw (16.0,1.4) to [short,*-] ++(0,-3.0) -- (15.4,-1.6);
\draw (13.2,-2.2) rectangle (15.4,-1.0);
\node[anchor=center] at (14.3,-1.6) {$\div\, 32$};
\draw (13.2,-1.6) -- (12.8,-1.6) -- (12.8,-2.45) -- (-1.2,-2.45) -- (-1.2,0.6) -- (0,0.6);
\node[anchor=south] at (6.0,-2.45) {\small CK\_FB};

% ---- bias and start-up ------------------------------------------------
% IBPSR_1U feeds the pump and the buffer; BIAS is its start-up helper
\draw (-0.8,-2.8) \portIn{IBPSR\_1U} (-0.8,-2.8) -- (10.6,-2.8);
\draw (5.0,-2.8) to [short,*-] (5.0,0);
\draw (8.3,-2.8) to [short,*-] (8.3,-3.4);
\draw (10.6,-2.8) -- (10.6,1.1);

\draw (2.2,-4.6) rectangle (4.0,-3.4);
\node[anchor=center] at (3.1,-4.0) {KICK};
\draw (7.4,-4.6) rectangle (9.2,-3.4);
\node[anchor=center] at (8.3,-4.0) {BIAS};

% the KICK pulse grabs the pump's LPFZ switch at start-up
\draw (4.0,-3.7) -- (6.2,-3.7) -- (6.2,0);
\node[anchor=south] at (5.0,-3.7) {\small KICK};
\draw (4.0,-4.3) -- (7.4,-4.3);
\node[anchor=north] at (5.7,-4.35) {\small $\overline{PWRUP}$};

% power-up
\draw (-0.8,-5.2) \portIn{PWRUP\_1V8} (-0.8,-5.2) -- (14.3,-5.2);
\draw (3.1,-5.2) to [short,*-] (3.1,-4.6);
\draw (5.6,-5.2) to [short,*-] (5.6,0);
\draw (14.3,-5.2) -- (14.3,-2.2);
\draw (13.7,0) to [short,-o] ++(0,-0.5);
\node[anchor=east] at (13.5,-0.5) {\small PWRUP\_1V8};

\end{circuitikz}
\end{document}
